Validating the \gls{DNP} group parameters is essential to ensure their practical utility.
Although the group parameters are verified with the literature, this does not necessarily mean they capture all physical phenomena of interest.
Validation identifies potential weaknesses in the approach by highlighting discrepancies between predicted and experimental results.

Validating the \gls{DNP} group parameters generated in this work is conducted in two stages, as discussed in Section \ref{sec:transimpact}.
Validation will be carried out in two stages, where the second stage builds upon the first.
The first stage validates the delayed neutron count rate formulation and group generation for a static problem.
The second stage validates the improved \gls{DNP} group parameters by comparing their reactivity feedback in a kinetic simulation to experimental results.

The first stage compares the delayed neutron count rate to the Keepin et al. experiment from 1957 \cite{keepin1957delayed}.
This requires simulated irradiation in Godiva in OpenMC and solving the k-eigenvalue problem to irradiate the sample.
The post-irradiation sample concentration in \gls{MoSDeN} is used to generate the \gls{DNP} group parameters.
The delayed neutron count rate, $\dot{n}_d$, and group parameters are compared with values from Keepin et al. to identify discrepancies.
Some differences are expected due to uncertainties in the data and experimental results, so validation is successful if the results align within uncertainty bounds.
It has already been shown that these results should have a percent difference of 9\% for thermal fission in $^{235}$U for \gls{JEFF}-3.1.1 data, so a similar percent difference can be expected for other datasets \cite{mathieu2012new}.

The second stage validates the improved \gls{DNP} group parameters.
This stage uses \gls{MSRE} pump start-up and coast-down experimental data to assess the accuracy of the group parameters in modeling reactivity changes from \glspl{DNP} \cite{prince1968zero}.
This stage requires modeling the \gls{MSRE} core and upper and lower plena in Moltres \cite{fratoni2020molten}.
The reactivity effects of the \glspl{DNP} are calculated for both transients using static and improved \gls{DNP} group parameters.
The improved group parameters are expected to predict reactivity more accurately as they account for more physical phenomena in the \gls{MSRE}.
If neither set of reactivity responses aligns with the experimental results, then there is an issue in the kinetics simulation, as Moltres has previously demonstrated the ability to model these transients \cite{reynolds2023analysis, park_advancements_2025}.
These transients have also been evaluated using other tools, providing additional comparison \cite{fei2021msre, delpech2003benchmark}.
If the static group parameters predict reactivity more accurately than the improved group parameters, this suggests an issue with the improved group parameter formulation, particularly if the first validation stage succeeded.
The issue must be resolved before proceeding, as the improved \gls{DNP} group parameters should provide higher accuracy in \gls{MSR} modeling.

%Potential challenges associated with the data are discussed in Section \ref{sec:dataeval}.